\chapter{Preliminaries}
\label{chap:preliminaries}

In this chapter we provide some background in optimization, list the standard assupmtions and definitions, and introduce the relevant probabilistic tools that will become central to our analysis in the sequel.

\section{Optimization}
The perfect optimization algorithm would take a function $f: \R^d \to \R$ and output a point $x^* \in \R^d$ satisfying
\begin{align*}
    x^* \in \operatorname*{argmin}_{x \in \R^d} f(x)
\end{align*}
However, since deciding whether a quartic polynomial $p(x_1, \ldots, x_d) \in \R[x_1, \ldots, x_d]$ is globally non-negative,
\begin{align*}
    &p(x) \geq 0 \qquad \text{for all } x \in \R^d
\end{align*}
is NP-hard~\cite{ahmadi2013nphardness}, it is proven in~\cite{anandkumar2016efficient} that even verifying whether a given point is a local minimum of a non-convex function is NP-hard in the worst case. So, the best we can hope for must be more restrictive. For general non-convex functions, the right definition ends up being
\begin{definition}
    A point $x \in \R^d$ is an $\ve$-first order stationary point if 
    \begin{align*}
        \mg{\grad f(x)} \leq \ve
    \end{align*}
\end{definition}
On the other hand, for convex $f$, we usually measure accuracy in terms of the function gap.
\begin{definition}
    Let $x^* \in \mathrm{argmin}_{x \in \R^d} f(x)$. A point $x \in \R^d$ is $\ve$-optimal if
    \begin{align*}
        f(x) - f(x^*) \leq \ve
    \end{align*}
\end{definition}
To reach a first-order stationary point, we need an idea of how much signal we get from the gradient versus noise. Recall that for any differentiable function $f(x)$,
\begin{align*}
    f(x+h) = f(x) + \langle \grad f(x), h \rangle + O(\mg{h}^2)
\end{align*}
The noise $O(\mg{h}^2)$ term is quantified by the $L$-smooth assumption
\begin{definition}
    A function $f \in \mathcal C^1(\R^d)$ is $L$-smooth if
    \begin{align*}
        \mg{\grad f(x) - \grad f(y)} \leq L \mg{x-y}
    \end{align*}
\end{definition}
Which leads to the descent lemma, the starting point to the analysis of any gradient-based algorithm.
\begin{lemma}[Descent Lemma]
    \label{lem:descent-lemma}
    If $f \in \mathcal C^1(\R^d)$ is $L$-smooth, then for any $h \in \R^d$,
    \begin{align*}
        f(x+h) \leq f(x) + \langle \grad f(x), h \rangle + \frac{L}{2}\mg{h}^2
    \end{align*}
\end{lemma}
\begin{proof}
    Writing $g(t) = f(x+th)$, we know that $\dv{t} g(t) = \langle \grad f(x+th), h \rangle$, and therefore
    \begin{align*}
        g(1) - g(0) = \int_0^1 \dv{t} g(t)\dd t = \int_{0}^1 \langle \grad f(x+th), h \rangle \dd t
    \end{align*}
    Therefore,
    \begin{align*}
        f(x+h) - f(x) - \langle \grad f(x),h\rangle = \int_{0}^1 \langle \grad f(x+th) - \grad f(x), h \rangle \dd t
    \end{align*}
    By the Cauchy-Schwarz inequality and the $L$-smooth assumption,
    \begin{align*}
        |\langle \grad f(x+th) - \grad f(x), h \rangle| &\leq \mg{\grad f(x+th) - \grad f(x)} \mg{h} \\
        &\leq tL\mg{h}^2
    \end{align*}
    Therefore,
    \begin{align*}
        |f(x+h) - f(x) - \langle \grad f(x),h\rangle| \leq L \mg{h}^2 \int_{0}^1 t \dd t = \frac{L}{2} \mg{h}^2
    \end{align*}
\end{proof}
The main optimization algorithm used in ML is gradient descent.
\begin{algorithm}[H]
\caption{Gradient Descent}
\label{alg:gradient-descent}
\begin{algorithmic}[1]
    \Require Objective function $f:\R^d\to\R$, smoothness constant $L>0$,
    initial point $x_0\in\R^d$, number of iterations $T$
    \For{$t=0,\ldots,T-1$}
        \State $x_{t+1} \gets x_t-\dfrac{1}{L}\grad f(x_t)$
    \EndFor
    \State \Return $x_T$
\end{algorithmic}
\end{algorithm}
By using the $L$-smoothness assumption, gradient descent converges to a first-order stationary point in $O(L (f(x) - f^*)/\ve^2)$ iterations where $f^* = f(x^*) = \mathrm{argmin}_{x \in \R^d} f(x)$. For convex functions, it demonstrates a linear convergence to an $\ve$-optimal point in $O(L\mg{x-x^*}^2/\ve)$ iterations. Gradient descent is the starting point for all algorithms discussed in the sequel.

\section{Zero-Order Oracles}
The focus of this work will be on zero-order optimization algorithms. Simply put, we drop the assumption that we can access gradients $\grad f(x)$. We can only evaluate the function (``oracle'') at points of our choosing $x \in \R^d, \; f(x)$. The simplest thing to do now is to calculate
\begin{align*}
    \widetilde{\grad} f(x) = \sum_{i=1}^d \frac{f(x+\sigma e_i) - f(x-\sigma e_i)}{2 \sigma} e_i 
\end{align*}
for any orthonormal basis $\set{e_1, \ldots, e_d}$ of $\R^d$. We have
\begin{align*}
    \mg{\widetilde{\grad} f(x) - \grad f(x)} \leq \frac{L\sigma \sqrt{d}}{2}
\end{align*}
This is a simple approach applied successfully, but is often of little practical use: frontier LLMs have $d$ in the trillions, and trillions of forward passes to calculate just one gradient is far worse than the usual backpropogation. 

For a quadratic $f(x) = \frac 12 x^\top H x$ with symmetric $H$, 
\begin{align*}
    \grad f(x) = Hx = \begin{bmatrix}
        \frac{f(x+e_1) - f(x-e_1)}{2} & \cdots & \frac{f(x+e_d) - f(x-e_d)}{2}
    \end{bmatrix}^\top
\end{align*}
Thus typically the gradient gives $d$ times more information than any single forward pass $f(x+e_i)$. This suggests that any zero-order optimization algorithm for convergence to a first-order stationary point for general nonconvex $f$ should take $d$ times more queries in the worst case. To the best of our knowledge this lower bound is currently open. Nevertheless, we are unaware of any method in the literature achieving a sublinear $d$ dependence for this problem.

Another way to use the oracle $f$ is to query random perturbations. For a random direction $p \sim \mathcal N(0,I_d)$, and search radius $\sigma$, we can look at $f(x+\sigma p)$. For $p_1, \ldots, p_r \overset{\mathrm{i.i.d.}}{\sim} \mathcal N(0,I_d)$ where $r$ is the population size, a Taylor expansion suggests the following estimator
\begin{align}
    \label{eq:estimator}
    \widetilde{\grad} f(x) = \frac 1r \sum_{i=1}^r \frac{f(x+\sigma p_i) - f(x-\sigma p_i)}{2\sigma} p_i.
\end{align}
In the sequel, this is the only estimator for the gradient we use. Upon defining the gaussian density 
\begin{align*}
    \cph_\sigma(z) = \frac{1}{(2\pi \sigma^2)^{d/2}} \exp\left(-\frac{\mg{z}^2}{2\sigma^2}\right)
\end{align*}
and the Gaussian smoothed objective
\begin{align*}
    f_\sigma(x) = (f * \cph_\sigma)(x) = \int_{\R^d} f(x+z)\cph_\sigma(z) \dd z = \mbb E_{p \sim \mathcal N(0,I_d)}[f(x+\sigma p)]
\end{align*}
under standard integrability assumptions
\begin{align*}
    \grad f_\sigma(x) = \frac{1}{\sigma} \E[f(x+\sigma p)p]
\end{align*}
which by symmetry of the Gaussian distribution yields,
\begin{align*}
    \grad f_\sigma(x) = \mbb E[\widetilde{\grad} f(x)]
\end{align*}
The advantage behind the central-difference version $\widetilde{\grad} f(x)$ and $\frac{1}{\sigma} f(x+\sigma p)p$ is that the first has good cancellation properties under a Taylor expansion
\begin{align*}
    f(x+\sigma p) - f(x-\sigma p) = 2\sigma \langle \grad f(x), p \rangle + O(\mg{\sigma p}^3).
\end{align*}
Since $\E[\widetilde{\grad} f(x)] = \grad f_\sigma(x)$, the bias of the estimator is precisely the error introduced by Gaussian smoothing. The high-dimensional analysis of this bias was investigated in~\cite{sarkar2025evolution}, where they proved 
\begin{align*}
    \mg{\E[\widetilde{\grad} f(x)] - \grad f(x)} = O_d(L \sigma \sqrt{d})
\end{align*}
Thus choosing $\sigma = o(d^{-1/2})$ ensures this bias vanishes at infinity. The intuition is that $f$ linearizes in high dimensions.
\section{The Low-Rank Trick}
In an ML model $\Phi_\theta(x)$, the parameters $\theta \in \R^{d \times m}$ are naturally thought of not as vectors but as matrices. While helpful that the estimator~\eqref{eq:estimator} lets us evaluate at only $r$ perturbations per iteration rather than the full $d \times m$, since $p \in \R^{d \times m}$ has all Gaussian entries, every entry is non-zero almost surely. Thus evaluating $\Phi_{\theta + \sigma p}(x)$ is still slow as we need to first add the two matrices $\theta + \sigma p$ before evaluating, with memory $O(rdm)$ just to store the perturbations.

The low rank trick, introduced in~\cite{sarkar2025evolution}, is to sample $a \sim \mathcal N(0,I_d)$, $b \sim \mathcal N(0, I_m)$, and exploit $(\theta + \sigma ab^\top)x = \theta x + \sigma a(b^\top x)$. This reduces the memory to $O(r(d+m))$, which makes zero-order algorithms practical.

However, for independent gaussian $Z,W \sim \mathcal N(0,1)$, the product $ZW$ is no longer subgaussian, it is sub-exponential, meaning it has fat tails. This could, \textit{a priori}, ruin the convergence of the algorithm, as $\E[\|\widetilde{\grad} f(x)\|^2]$ could grow too large. This is the central question of Chapter~\ref{chap:fosp-eggroll}. To heed our analysis, for $A,B \in \R^{d \times m}$, we define the Frobenius inner product and the Frobenius norm,
\begin{align*}
    \langle A, B \rangle_F = \Tr(AB^\top) = \sum_{i,j} A_{i,j} B_{i,j}, \quad \mg{A}_F = \sqrt{\sum_{i,j} A_{i,j}^2}
\end{align*}
We define $L$-smoothness in the same way as before using this norm. 

The Gaussian distribution is rotationally invariant, meaning for $p \sim \mathcal N(0,I)$, for any rotation matrix $R$ satisfying $RR^\top = R^\top R = I$, we have $Rp \sim \mathcal N(0,I)$. Thus, we search all directions uniformly. However, $ab^\top$ does not have this property, as it is rank 1, we could design a rotation matrix which rotates it into $M$ which is full rank. This raises the question, is rotational invariance necessary? What if we miss the direction of the gradient?
\begin{definition}
    The random vector $X \in \R^d$ is isotropic if
    \begin{align*}
        \E[XX^\top] = I
    \end{align*}
    this is equivalent to saying, for any $u \in \R^d$,
    \begin{align*}
        \E[\langle X, u \rangle^2] = \E[u^\top XX^\top u] = u^\top \E[XX^\top]u = \mg{u}^2
    \end{align*}
\end{definition}
While $ab^\top$ is not rotationally invariant, it is isotropic:
\begin{lemma}
    For $a \sim \mathcal N(0,I_d), b \sim \mathcal N(0, I_m)$ independent, for any $U \in \R^{d \times m}$,
    \begin{align*}
        \E[\langle ab^\top, U \rangle_F^2] = \mg{U}_F^2
    \end{align*}
\end{lemma}
\begin{proof}
    We have
    \begin{align*}
        \E[\langle ab^\top, U \rangle_F^2] = \E[(\Tr(U b a^\top))^2]
    \end{align*}
    By cyclic property of trace, we know that $\Tr(ABC)=\Tr(CAB)$, so the above equals
    \begin{align*}
        \E[(a^\top U b)^2] = \E[b^\top U^\top a a^\top U b]
    \end{align*}
    By conditioning on $b$, using that $a$ and $b$ are independent, and that $a \sim \mathcal N(0,I)$,
    \begin{align*}
        \E[b^\top \E[U^\top a a^\top U \mid b] b] = \E[b^\top U^\top \E[aa^\top] U b] = \E[b^\top U^\top U b]
    \end{align*}
    This equals
    \begin{align*}
        \E[\Tr(U^\top U bb^\top)] = \Tr(U^\top U \E[bb^\top]) = \Tr(U^\top U) = \mg{U}_F^2.
    \end{align*}
\end{proof}
It will turn out that isotropy is exactly the condition we need for convergence to a first-order stationary point. 

\section{Probability Tools}
Here we give a short background on measure-theoretic probability. We define a $\sigma$-algebra
\begin{definition}
    For a state space $\Omega$, $\mathcal F \subset 2^\Omega$ is a $\sigma$-algebra if:
    \begin{enumerate}[nosep]
        \item $\Omega \in \mathcal F$
        \item $\mathcal F$ is closed under complements: if $A \in \mathcal F$, then $A^c \in \mathcal F$
        \item $\mathcal F$ is closed under countable unions: if $A_1, A_2, \ldots \in \mathcal F$, then $A = A_1 \cup A_2 \cup \cdots \in \mathcal F$. 
    \end{enumerate}
\end{definition}
By combining the last two properties $\mathcal F$ is also closed under countable intersections. $\mathcal F$ is the collection of \textit{measurable sets}, or in probability, \textit{events}. 
\begin{definition}
    $\mbb P: \mathcal F \to [0,1]$ is a probability measure if $\P(\emptyset) = 0$, $\P(\Omega) = 1$, and
    \begin{align*}
        \P(\bigcup_{i=0}^{\infty} E_i) = \sum_{i=0}^{\infty} \P(E_i)
    \end{align*}
\end{definition}
The collection $(\Omega, \mathcal F, \mbb P)$ is called a probability space. 

In stochastic optimization, we often want to talk about the information available at iteration $t$. Filtrations give us a natural way to do this. 
\begin{definition}
    $\set{\mathcal F_t}_{t=0}^\infty$ is a filtration on $(\Omega, \mathcal F, \mbb P)$ if each $\mathcal F_t \subset \mathcal F$ is $\sigma$-algebra, and
    \begin{align*}
        \mathcal F_0 \subset \mathcal F_1 \subset \mathcal F_2 \subset \cdots
    \end{align*}
\end{definition}
For Chapters~\ref{chap:zo-mirror-descent} and~\ref{chap:linear-mixing}, at each iteration $t$ we sample $p_{t,1}, \ldots, p_{t,r} \sim \mathcal N(0,I_d)$ to evaluate $\widetilde \grad f(x_t)$. Thus we define the following filtration
\begin{definition}
    The zero-order filtration $\set{\mathcal F_t}_{t=0}^\infty$ is defined as
    \begin{align*}
        \mathcal F_{t-1}
    =
    \sigma\left(
        x_0,\ldots,x_t,
        \{p_{\tau,i}:0\leq \tau\leq t-1,\ 1\leq i\leq r\}
    \right).
    \end{align*}
    which encapsulates all information available at time $t$. 
\end{definition}
Central to our analysis are concentration inequalities. Recall Hoeffding's inequality:
\begin{lemma}[Hoeffding's inequality]
    Let $X_1,\ldots,X_n$ be independent random variables satisfying
    $a_i \leq X_i \leq b_i$ almost surely for every $1\leq i\leq n$.
    Then, for every $t>0$,
    \[
        \Pr\left(
            \left|
                \sum_{i=1}^n \big(X_i - \E[X_i]\big)
            \right|
            \geq t
        \right)
        \leq
        2\exp\left(
            -\frac{2t^2}
            {\sum_{i=1}^n (b_i-a_i)^2}
        \right).
    \]
\end{lemma}
By re-arranging, we can get
\begin{lemma}
    Let $X_1,\ldots,X_n$ be independent random variables satisfying
    $a_i \leq X_i \leq b_i$ almost surely for every $1\leq i\leq n$. Then with probability at least $1-\delta$, we have
    \begin{align*}
        \left|\sum_{i=1}^n \big(X_i - \E[X_i]\big)\right| \leq \sqrt{\frac 12\sum_{i=1}^n (b_i-a_i)^2 \log(2/\delta)}
    \end{align*}
\end{lemma}
Using $a_i \leq X_i \leq b_i$, $\Var(X_i) = \E[(X_i-\E(X_i))^2] \leq (b_i-a_i)^2$. A bernoulli random variable which is $a_i$ and $b_i$ with probability 1/2 has variance $(b_i-a_i)^2/4$, so this is tight up to constant factors. Since $X_i$ are independent, $\Var(X_1 + \cdots + X_n) = \Var(X_1) + \cdots + \Var(X_n) = \sum_{i=1}^n (b_i-a_i)^2$. Hoeffding's inequality says that, when the $a_i \leq X_i \leq b_i$, the absolute value of the centered sum is not typically that much larger than its expectation. 

However, this property is not limited to just $a_i \leq X_i \leq b_i$. It carries over to martingale difference sequences. 
\begin{definition}[Martingale Difference Sequence]
    $\set{\xi_t}_{t\geq 1}$ is a martingale difference sequence wrt filtration $\mathcal F_t$ if $\xi_t$ is $\mathcal F_t$-measurable and 
    \begin{align*}
        \E[\xi_t \mid \mathcal F_{t-1}] = 0
    \end{align*}
\end{definition}
The main intuition is that in a martingale difference sequence, the variance of the full sum is the sum of the variance of each $\xi_t$, which is vastly superior than including all the cross covariance terms $\Cov(\xi_i, \xi_j)$. This improves our bound substantially. 
\begin{lemma}
    Define $S_T = \sum_{t=1}^T \xi_t$. Then
    \begin{align*}
        \E[S_T^2] = \sum_{t=1}^T \E[\xi_t^2]
    \end{align*}
\end{lemma}
\begin{proof}
    We have
    \begin{align*}
        \E[S_T^2] = \sum_{t=1}^T \E[\xi_t^2] + 2\E[\sum_{i < j} \xi_i \xi_j]
    \end{align*}
    Now,
    \begin{align*}
        \E[\sum_{i < j} \xi_i \xi_j] = \sum_{i < j} \E[\E[\xi_i \xi_j \given \mathcal F_{j-1}]] = \sum_{i < j} \E[\xi_i \E[ \xi_j \given \mathcal F_{j-1}]] = 0.
    \end{align*}
\end{proof}

For us, these martingale terms $\xi_i$ will always be subgaussian.
\begin{definition}
    The random variable $X \in \R$ is $\sigma$-subgaussian if
    \begin{align*}
        \E[\exp(\frac{X^2}{3\sigma^2})] \leq 2
    \end{align*}
\end{definition}
Under this definition $\mathcal N(0,1)$ is 1-subgaussian (the constant 3 is not standard). We have the following lemma which vastly generalizes Hoeffding's inequality and will help us substantially in the sequel,
\begin{lemma}[Adapted Sub-Gaussian Hoeffding's inequality]
    Let $\{\mathcal F_t\}_{t=0}^T$ be a filtration, and let $X_t$ be
    $\mathcal F_t$-measurable. Suppose that, conditional on
    $\mathcal F_{t-1}$, $X_t$ is $\sigma_t$-subgaussian and mean zero.
    Then, for any $\delta\in(0,1)$, with probability at least $1-\delta$,
    \[
        \left|\sum_{t=1}^T X_t\right|
        \leq
        c\sqrt{
            \log\left(\frac{2}{\delta}\right)
            \sum_{t=1}^T \sigma_t^2
        },
    \]
    where $c>0$ is an absolute constant.
\end{lemma}

Convergence of stochastic gradient methods are standard in the literature, see for example~\cite{ghadimi2013stochastic,nesterov2017random,balasubramanian2018zeroth}. Concentration inequalities will allow us to upgrade that convergence in expectation to convergence guarantees that hold with high probability.

\section{Notation}
$d$ and $m$ will always be dimension constants, $r$ will always be the population size, $T$ will always be the iteration horizon, $L$ will always be the smoothness constant, $\sigma$ will always be the smoothing radius, $\ve$ will always be the desired accuracy, and $\delta$ will always be the failure probability. We use $O(\cdot), \Omega(\cdot)$, and $\Theta(\cdot)$ in the standard sense. The notation $\widetilde O(\cdot)$ hides polylogarithmic factors in the problem parameters, including $d$, $L$, $T$, $f(x_0)-f^*$, $\ve$ and failure probability $\delta$. Throughout the sequel, $\chi$ always denotes a polylogarithmic factor.
